\begin{helper}
For every real number \(0<\eta<1/2\), there is a real number
\(\varepsilon>0\) such that, for every integer \(n_0\), there is an integer
\(k_0\) with the following property.  For every integer \(k\ge k_0\), there is
an integer \(N\) such that
\[
n_0\le N,
\]
\[
\left(\frac12-\eta\right)\noderef{RamseyScale}(k)\le (N:\mathbb R),
\]
and
\[
(1+\varepsilon)\sqrt{(N:\mathbb R)\log(N:\mathbb R)}<(k:\mathbb R).
\]
\end{helper}
\begin{proof}
Fix \(0<\eta<1/2\).  Put
\[
c=\frac12-\frac{\eta}{2}
\qquad\text{and}\qquad
\delta=c-\left(\frac12-\eta\right)=\frac{\eta}{2}.
\]
Then \(0<c<1/2\), \(\frac12-\eta<c\), and \(\delta>0\).  Since \(2c<1\), choose
\(\varepsilon>0\) so small that
\[
2c(1+\varepsilon)^2<1. \tag{1}
\]
This is possible because the function \(x\mapsto 2c(1+x)^2\) is continuous at
\(0\), and its value there is \(2c<1\).

Now fix \(n_0\).  Choose \(k_0\) so large that for every \(k\ge k_0\),
\[
k\ge3,\qquad \log k>0,\qquad \log k>c,
\]
\[
c\,\frac{k^2}{\log k}\ge n_0+1,\qquad
c\,\frac{k^2}{\log k}\ge4,\qquad
\delta\,\frac{k^2}{\log k}\ge1. \tag{2}
\]
Here is the threshold justification.  Since \(\log k\to\infty\), the
conditions \(\log k>0\) and \(\log k>c\) hold for all sufficiently large
\(k\).  For all sufficiently large \(k\) we also have \(0<\log k\le k\); hence
\[
\frac{k^2}{\log k}\ge k.
\]
Thus \(k^2/\log k\) tends to infinity.  Therefore the three lower bounds
involving \(k^2/\log k\) also hold eventually.  Taking \(k_0\) to be the maximum
of these finitely many thresholds gives (2) simultaneously.

Fix \(k\ge k_0\), and set
\[
X=c\,\frac{k^2}{\log k},\qquad N=\lfloor X\rfloor.
\]
By (2), \(X\ge0\).  The floor inequalities for a nonnegative real number give
\[
(N:\mathbb R)\le X
\qquad\text{and}\qquad
X-1\le (N:\mathbb R). \tag{3}
\]
The last inequality in (2) gives
\[
1\le \delta\,\frac{k^2}{\log k}.
\]
Subtracting this bound for \(1\) from \(X=c\,k^2/\log k\), we get
\[
X-1\ge
\left(c-\delta\right)\frac{k^2}{\log k}
=\left(\frac12-\eta\right)\frac{k^2}{\log k}. \tag{4}
\]
Combining (3), (4), and the definition of \(\noderef{RamseyScale}\), we obtain
\[
\left(\frac12-\eta\right)\noderef{RamseyScale}(k)\le (N:\mathbb R). \tag{5}
\]
The first middle inequality in (2) gives
\[
(N:\mathbb R)\ge X-1\ge n_0,
\]
so \(n_0\le N\).  The second middle inequality in (2) gives
\[
(N:\mathbb R)\ge X-1\ge3,
\]
so \(N\ge3\).

It remains to prove the square-root estimate.  Since \(N\ge3\), \(0<N\) and
\(\log N>0\).  From \((N:\mathbb R)\le X=c k^2/\log k\) and \(\log k>c>0\),
\[
(N:\mathbb R)\le c\,\frac{k^2}{\log k}<k^2. \tag{6}
\]
Thus \(0<N\le k^2\).  Monotonicity of the logarithm on positive reals gives
\[
\log N\le \log(k^2)=2\log k. \tag{7}
\]
Multiplying the bounds (3) and (7), using \(\log k>0\), yields
\[
N\log N\le
\left(c\,\frac{k^2}{\log k}\right)(2\log k)=2c\,k^2. \tag{8}
\]
Multiplying (8) by the positive number \((1+\varepsilon)^2\) and using (1),
\[
(1+\varepsilon)^2N\log N<k^2. \tag{9}
\]
Both sides are nonnegative and \(k>0\).  Taking square roots and using
\[
\sqrt{(1+\varepsilon)^2N\log N}
=(1+\varepsilon)\sqrt{N\log N}
\]
gives
\[
(1+\varepsilon)\sqrt{N\log N}<k.
\]
Together with (5) and \(n_0\le N\), this is the required integer \(N\).
\end{proof}
